\section{Simultaneous recoding of formula codes and variable labels}
A future change of constructor encoding need not create a mismatch between
an axiom's property code and its bound-variable label. The numerical
calculus uses every natural number in both roles. Recode both roles with the
same recursively defined permutation. This observation is useful at the
calculus level; no smaller literal machine is obtained here.

\paragraph{An explicit alternative.}
Keep the transposed quantifier convention and define
\[
\begin{aligned}
 I'(a,b)&=2p(a,b),&M'(a,b)&=6p(a,b)+1,\\
 U'(a,B)&=6p(B,a)+3,&E'(a,b)&=6p(a,b)+5.
\end{aligned}
\]
Even numbers are implication codes. Odd residues $1,3,5$ modulo six are
membership, universal-quantifier, and equality codes respectively. These
ranges partition $\mathbb N$. Code zero remains the truth alias
$I'(0,0)=0$, and code one remains $M'(0,0)=1$. Every component of a nonzero
constructor is strictly smaller than its code, including both components
of atoms, which are variable labels.

Let $C_r(a,b)=4p(a,b)+r$ be the old raw constructors, where the raw
quantifier components are body first and binder second. Let $C'_0=I'$,
$C'_1=M'$, $C'_2(a,b)=U'(b,a)$, and $C'_3=E'$. Define
\[
 f(0)=0,\qquad
 f(C_r(a,b))=C'_r(f(a),f(b))\quad(C_r(a,b)>0).
\]
Define $g$ by the same recursion in the reverse direction, using the unique
new tag and pair decomposition. The zero clause is necessary: blindly
recursing on $C_0(0,0)$ would not descend.

\paragraph{Constructor-algebra lemma.}
The maps $f,g$ are primitive recursive inverse bijections. They fix zero and
one and satisfy, for every label or formula code,
\[
\begin{aligned}
 f(I(a,b))&=I'(f(a),f(b)),&
f(M(a,b))&=M'(f(a),f(b)),\\
 f(U(a,B))&=U'(f(a),f(B)),&
f(E(a,b))&=E'(f(a),f(b)).
\end{aligned}
\]
To prove this, compute a course-of-values table in increasing order of the
input code. Every recursive argument is smaller. Explicit primitive
recursive growth bounds are available: if all earlier values are at most
$m$, each new forward value is at most $12m^2+12m+5$, and each new reverse
value at most $8m^2+8m+3$. Iterating these bounds supplies a primitive
recursive bound for each finite table. Induction on the old code proves
$g(f(n))=n$, and induction on the new code proves $f(g(n))=n$, by unique
constructor decomposition. The homomorphism equations follow from the
recursion, including the exceptional $I(0,0)$ case. These arguments use only
ordinary arithmetic induction and admit uniform formalization in PA.

This is not the translation that fixes every numerical variable label while
changing only compound formula codes. A variable $x_a$ becomes $x_{f(a)}$.
Consequently a gate using $c$ as both property code and binder becomes a
gate using $f(c)$ in both roles. The guards $c\ne a,b$ are preserved by
injectivity. For example, if
\[
 F(a,b,c)=N\bigl(I(\Delta_c(a,b),I(V(a,b,c),Q(a,b,c)))\bigr),
\]
then exact constructor homomorphism gives
\[
 f(F(a,b,c))=F'(f(a),f(b),f(c)).
\]
The same statement applies to the diagonal clause $\Gamma(c)=I(c,U(c,c))$
and to its separated or bundled gates. No independently asserted
alpha-conversion or freshness shortcut is needed for this identity.

\paragraph{The numerical order restriction must be handled separately.}
The map is not monotone:
\[
 f(2)=3,\qquad f(3)=5,\qquad f(4)=2.
\]
Thus the hypothesis $u\le a$ of a restricted $\rho_8$ instance does not
transport to $f(u)\le f(a)$. The correct proof translation first uses the
unrestricted implicational schema already derived from the restricted one.
The construction in the imported
\texttt{M6/Logic/RestrictedImplication.lean} uses only reflexivity of the
parameter order and domination of both arguments by the implication code.
The new constructor $I'(a,b)=2p(a,b)$ has these same domination properties.
Hence the same finite derivation recipes yield unrestricted $\rho'_8$.
Map an old restricted row to that derived unrestricted new row; use the
analogous old theorem in the reverse direction.

For the six quantifier/equality heads and the gate, the image is a literal
instance of the corresponding new schema. Modus ponens and generalization
transport by constructor homomorphism. Induction on coded derivations then
gives primitive recursive proof transformations in both directions,
preserving their recoded conclusions. Since $f(1)=g(1)=1$, the resulting
abstract numerical calculi have PA-provably equivalent refutation
predicates. This is a written calculus-level argument and a route for a
future formal transport module; it is not a claim that a modified register
program, its command compiler, or its literal tape simulator has been
verified.

\paragraph{Physical tradeoff.}
The alternative makes an implication constructor multiply its Cantor pair by
two, and an MP recognizer need only check evenness before unpairing. It also
makes each nonimplicational constructor multiply by six, which introduces
additional arithmetic and changes the emitter's sharing patterns. The
present program shares its constructor and decoding fragments extensively;
source instruction savings cannot simply be subtracted from its working
state count. In particular, replacing only the factor four in the existing
cons routine would give incorrect membership, quantifier, equality, and
proof-command behavior. A complete generator, decoder, command simulation,
startup check, and allocation would be required before reporting a state
count for this recoding.

The accompanying \texttt{recoding.py} checks both inverse identities on
10,000 inputs, 1,764 constructor identities, and 64 exact gate identities.
These finite checks verify the transcription of the definitions; the
inductive argument above establishes the general algebraic claims.
